\documentclass[10pt,a4paper]{amsart}
\usepackage[margin=19mm]{geometry}
\usepackage{amsmath,amssymb,mathtools}
\usepackage[scr=rsfs,cal=euler]{mathalfa}
\usepackage{xfrac}
\usepackage[unicode,hidelinks,bookmarks=false]{hyperref}
\newcommand{\ie}{i.e.\ }
\newcommand{\cf}{cf.\ }
\newcommand{\resp}{respectively\ }
\newcommand{\Subsection}{\subsection}
\providecommand{\nobreakdash}{\nobreak\hbox{-}}
\setlength{\emergencystretch}{2em}
\multlinegap0pt
\usepackage{amssymb}
\usepackage[normalem]{ulem}
\usepackage{xcolor}
\usepackage[shortlabels]{enumitem}
\newtheorem{theorem}{Theorem}
\newtheorem{proposition}{Proposition}
\newtheorem{corollary}{Corollary}
\def\N{\mathbb{N}}
\def\R{\mathbb{R}}
\def\Z{\mathbb{Z}}
\renewcommand{\phi}{\varphi}
\title{Smooth linearization of nonautonomous dynamics under general dichotomic behaviour}
\author{Lucas Backes, Davor Dragicevic, Wenmeng Zhang}
\date{Source: arXiv:2605.04184v1 (CC BY 4.0)}
\begin{document}
\maketitle

\noindent\emph{Source and licence.} Smooth linearization of nonautonomous dynamics under general dichotomic behaviour. Version arXiv:2605.04184v1 (CC BY 4.0), distributed by arXiv under the Creative Commons Attribution 4.0 International licence, http://creativecommons.org/licenses/by/4.0/, as recorded in the arXiv licence metadata for this exact version and shown on the abstract page https://arxiv.org/abs/2605.04184v1. Full licence notice: \texttt{licenses/arxiv-2605.04184-CC-BY-4.0.txt}.\par\smallskip

\noindent\textit{Editorial scope.} Counted unit: the article's theorem pol (source0.tex lines 1117--1148). The article's complete original proof of it is delivered with it (source0.tex lines 1150--1327, one \texttt{proof} environment of 178 lines). No mathematics is written, repaired or re-derived: the statement and the proof are the article's own lines, in the article's own notation, and no retained line is substituted or re-flowed.  Retained uncounted as the source-local context the proof needs: lines 293--295; lines 367--369; lines 394--396; lines 398--404; lines 530--558; lines 536--538; lines 540--542; lines 545--547; lines 550--552; lines 554--556; lines 963--965; lines 968--970; lines 995--997; lines 1004--1018; lines 1006--1008; lines 1014--1016; lines 1108--1114.  Nothing was omitted from any retained span, so the package carries no omission record.  No retained line contains a citation, so the wrapper carries no bibliography and no bibliography entry.  No retained line contains a figure environment, an image-inclusion command or drawing code, so no image is delivered and the package declares no asset dependency.  Editorial formatting only, in addition: an AMS \texttt{amsart} preamble that restates the article's own declarations and macros used by the retained lines, namely the environments \texttt{theorem}, \texttt{proposition}, \texttt{corollary} on separate counters, together with the standard packages those lines need.  The retained environments print in this document as Theorem 1, Proposition 1, Corollary 1 in order of appearance, whereas the article prints the same items with the numbering of its own full document.  The complete native source of the article as distributed in the arXiv e-print for this exact version is bundled unchanged as \texttt{sources/199840/source0.tex}.
\begin{equation}\label{eq: grbound}
	\frac{\mu_{n+1}}{\mu_n}\le \theta, \quad \text{for all} \quad n \in \N,
\end{equation} 

\begin{equation}\label{EDB-mu}
\Sigma_{\mu D, \mathbb A}=\bigcup_{i=1}^r [a_i, b_i].
\end{equation}

\begin{equation}\label{sp1}
a_1\le b_1 <\ldots <a_k\le b_k <0<a_{k+1}\le b_{k+1}<\ldots <a_r\le b_r,
\end{equation}

\begin{equation}\label{sp2}
\begin{cases}
a_{k+1}-b_k>\max \{b_r, -a_1\},\\
b_i-a_i\le -b_k \text{ for } 1\le i\le k,\\
b_i-a_i\le a_{k+1} \text{ for } k+1\le i \le r.\\
\end{cases}
\end{equation}

\begin{theorem}\label{theo: main linearization discrete}
Let $\mu=(\mu_n)_{n\in \Z^+}$ be a growth rate that satisfies~\eqref{eq: grbound} for some $\theta \ge 1$ and $\mathbb A=(A_n)_{n\in \mathbb N}$ be a sequence of invertible linear operators on $\mathbb R^d$ that admits a strong $\mu$-dichotomy. Furthermore, suppose $\Sigma_{\mu D, \mathbb A}$ has the form \eqref{EDB-mu}
with $1\le r\le d$ and that the numbers $a_i$ and $b_i$ satisfy \eqref{sp1} and \eqref{sp2}.
Finally, let $(g_n)_{n\in \N}$ be a sequence of $C^1$ maps $g_n\colon \R^d \to \R^d$ satisfying the following:
\begin{itemize}
\item for every $n\in \N$,
\begin{equation}\label{eq: gn 0 and Dgn 0}
g_n(0)=0 \;\text{ and } \; Dg_n(0)=0;
\end{equation}
\item there exists $c>0$ such that 
\begin{equation}\label{eq: bound deriv g}
\|Dg_n(x)\| \le c\frac{\mu_n'}{\mu_n}, \quad \text{for $x\in \R^d$ and $n\in \N$,}
\end{equation}
where $\mu_n'=\mu_{n+1}-\mu_n$;
\item there exists $M>0$ such that 
\begin{equation}\label{4-4-3}
\|Dg_n(x)-Dg_n(y)\| \le M\frac{\mu_n'}{\mu_n}\|x-y\|, \quad \text{for $x, y\in \R^d$ and $n\in \N$.}
\end{equation}
\end{itemize}
Then, provided that $c$ is sufficiently small, there exists a sequence $(\psi_n)_{n\in \N}$ of $C^1$-diffeomorphisms $\psi_n\colon \R^d \to \R^d$ such that
\begin{equation}\label{lincond}
\psi_{n+1}\circ (A_n+g_n)=A_n  \circ \psi_n, \; \text{ for } n\in \N,
\end{equation}
and there exist $T, \rho>0$ so that 
\begin{equation}\label{hn-1}
\|D\psi_n(x)\| \le T \quad \text{and} \quad \|D\psi_n^{-1}(x)\| \le T,
\end{equation}
for $n\in \N$ and $x\in \R^d$ with $\|x\|\le \rho$.
\end{theorem}

\begin{equation}
g_n(0)=0 \;\text{ and } \; Dg_n(0)=0;
\end{equation}

\begin{equation}
\|Dg_n(x)\| \le c\frac{\mu_n'}{\mu_n}, \quad \text{for $x\in \R^d$ and $n\in \N$,}
\end{equation}

\begin{equation}
\|Dg_n(x)-Dg_n(y)\| \le M\frac{\mu_n'}{\mu_n}\|x-y\|, \quad \text{for $x, y\in \R^d$ and $n\in \N$.}
\end{equation}

\begin{equation}
\psi_{n+1}\circ (A_n+g_n)=A_n  \circ \psi_n, \; \text{ for } n\in \N,
\end{equation}

\begin{equation}
\|D\psi_n(x)\| \le T \quad \text{and} \quad \|D\psi_n^{-1}(x)\| \le T,
\end{equation}

\begin{equation}\label{LDE}
x'=A(t)x, \quad t\ge 1,
\end{equation}

\begin{equation}\label{NDE}
x'=A(t)x+f(t,x), \quad t\ge 1.
\end{equation}

\begin{equation}\label{bnd}
\| T(t,s)\| \le K\left(\frac{\mu(t)}{\mu(s)} \right)^a \quad \text{and} \quad \| T(s,t)\| \le K \left(\frac{\mu(t)}{\mu(s)}\right)^a,
\end{equation}

\begin{proposition} \label{prop: equiv pol dich contXdisc}
Assume that there exist $K,a>0$ such that the evolution family $T(t,s)$ of \eqref{LDE} satisfies \eqref{bnd}. Moreover, suppose that the growth rate $\mu$ satisfies
\begin{equation}\label{eq: grbound-cont}
	\frac{\mu(t+1)}{\mu (t)}\le \theta, \quad \text{for all} \quad t\geq 1,
\end{equation} 
for some $\theta\ge 1$
Then, the following properties are equivalent:
\begin{itemize}
\item \eqref{LDE} admits a strong $\mu$-dichotomy;
\item the sequence $\mathbb A =(A_n)_{n\in \N}$ admits a strong $\mu$-dichotomy, where
\begin{equation}\label{AN}
A_n=T(n+1, n), \quad n\in \N.
\end{equation}
\end{itemize}
\end{proposition}

\begin{equation}
	\frac{\mu(t+1)}{\mu (t)}\le \theta, \quad \text{for all} \quad t\geq 1,
\end{equation} 

\begin{equation}
A_n=T(n+1, n), \quad n\in \N.
\end{equation}

\begin{corollary}\label{cor: equal spectrum cont}
Assume that there exist $K, a>0$ such that~\eqref{bnd} holds and that $\mu$ satisfies \eqref{eq: grbound-cont} for some $\theta\geq 1$. Then
\[
\Sigma_{\mu D, A(\cdot)}=\Sigma_{\mu D, \mathbb A},
\]
where the sequence $\mathbb A=(A_n)_{n\in \N}$ is given by~\eqref{AN}. In particular, $\Sigma_{\mu D, A(\cdot)}$ is given by~\eqref{EDB-mu} with $1\leq r\le d$ and $a_1\leq b_1<a_2\leq \ldots < a_r\leq b_r$.
\end{corollary}

\begin{theorem}\label{theo: linarization continuous time pol}
Let $\mu$ be a differentiable growth rate satisfying \eqref{eq: grbound-cont} for some $\theta\geq 1$ and suppose that \eqref{LDE} admits a strong $\mu$-dichotomy. Assume that $\Sigma_{\mu D, A(\cdot)}$ is given by~\eqref{EDB-mu}, where $a_i$ and $b_i$ satisfy~\eqref{sp1} and~\eqref{sp2}.
 Furthermore, suppose that $f\colon [1, \infty)\times \R^d \to \R^d$ is a $C^1$ map satisfying the following conditions:
\begin{itemize}
\item for $t\ge 1$,
\begin{equation}\label{1950}
f(t,0)=0 \quad \text{and} \quad D_x f(t,0)=0;
\end{equation}
\item there exists $\eta >0$ such that 
\begin{equation}\label{1951}
\| D_x f(t,x)\| \le \eta \frac{\mu'(t)}{\mu(t)}, \quad \text{for $x\in \R^d$};
\end{equation}
\item there exists $L>0$ such that 
\begin{equation}\label{1952}
\|D_x f(t,x)-D_x f(t, y)\| \le L\frac{\mu'(t)}{\mu(t)} \|x-y\|, \quad \text{for $x, y\in \R^d$}.
\end{equation}
\end{itemize}
Then, provided that $\eta$ is sufficiently small,  there exist $C^1$ maps $H, G\colon [1, \infty) \times \R^d \to \R^d$ with the following properties:
\begin{itemize}
\item[i)] if $t\mapsto x(t)$ is a solution of~\eqref{NDE}, then $t\mapsto H(t, x(t))$ is a solution of~\eqref{LDE};
\item[ii)] if $t\mapsto x(t)$ is a solution of~\eqref{LDE}, then $t\mapsto G(t, x(t))$ is a solution of~\eqref{NDE};
\item[iii)] for $t\ge 1$ and $x\in \R^d$, 
\[
H(t, G(t,x))=x \quad \text{and} \quad G(t, H(t,x))=x;
\]
\item[iv)] there exist $R, \zeta>0$ such that 
\[
\|D_xH(t,x)\|\leq R \; \text{ and }\; \|D_xG(t,x)\| \leq R , 
\]
for every $t\geq 1$ and $x\in \R^d$ that satisfies $\|x\|\leq \zeta$.
\end{itemize}
\end{theorem}

\begin{proof} The general idea of the proof consists of using the discretization of \eqref{LDE} given by \eqref{AN} to obtain Theorem \ref{theo: linarization continuous time pol} as a consequence of Theorem \ref{theo: main linearization discrete}.


Let $(A_n)_{n\in \N}$ be the sequence of operators on $\R^d$ given by~\eqref{AN}. Combining our assumptions with Proposition \ref{prop: equiv pol dich contXdisc} and Corollary \ref{cor: equal spectrum cont}, we easily see that $(A_n)_{n\in \N}$ satisfies the hypotheses of Theorem \ref{theo: main linearization discrete}. Let $t\mapsto \varphi (t, t_0; x_0)$ denote the solution of~\eqref{NDE} satisfying $x(t_0)=x_0$. Hence, by the variation of constants formula, we have that 
\begin{align}\label{def-phi}
\phi(t, n; x)=T(t,n)x+\int_n^t T(t,r)f(r, \phi(r, n;x))\, dr.
\end{align}
For $n\in \N$ and $x\in \R^d$, set 
\begin{equation}\label{eq: def gn cont case}
g_n(x)=\phi(n+1, n;x)-A_n x=\int_n^{n+1}T(n+1, r)f(r, \phi(r,n;x))\, dr.
\end{equation}
We will now verify that this sequence of $C^1$ maps $g_n:\mathbb{R}^d\to \mathbb{R}^d$, $n\in \N$ satisfies all the assumptions of Theorem~\ref{theo: main linearization discrete}.

It follows from~\eqref{1950} that  $\phi(t, t_0;0)=0$, and thus $g_n(0)=0$ for $n\in \N$. Moreover, 
\[
\begin{split}
Dg_n(0)
&=\int_n^{n+1}T(n+1, r)D_xf(r, 0)D_x \phi(r,n;0)\, dr =0,
\end{split}
\]
for $n\in \N$. We conclude that~\eqref{eq: gn 0 and Dgn 0} is true.

 Observe that 
\begin{equation}\label{dphi}
D_x\phi(t, n; x)=T(t,n)+\int_n^t T(t,r)D_xf(r, \phi(r, n;x))D_x\phi(r, n;x)\, dr
\end{equation}
for $t\ge n$ and $x\in \R^d$.  In particular, from~\eqref{bnd}, \eqref{eq: grbound-cont} and~\eqref{1951} we obtain that for every $t\in [n, n+1]$ and $x\in \R^d$,
\[
\begin{split}
\lVert D_x\phi(t, n; x) \rVert&\le K\bigg (\frac{\mu (t)}{\mu (n)} \bigg )^a  + \int_n^t K \bigg (\frac{\mu(t)}{\mu(r)} \bigg )^a \eta \frac{\mu'(r)}{\mu(r)} \lVert D_x\phi(r, n;x)\rVert\, dr
\\
&\le K\theta^a  +\int_n^t K\eta  \theta^a  \frac{\mu'(r)}{\mu(r)} \lVert D_x\phi(r, n;x)\rVert\, dr.
\end{split}
\]
 Hence, from Gronwall's lemma it follows that 
$\lVert D_x\phi(t, n; x) \rVert \le K \theta^a  e^{\int_n^t K\eta  \theta^a  \frac{\mu'(r)}{\mu(r)}\, dr}$. Thus, since (using \eqref{eq: grbound-cont})
\begin{equation}\label{eq: aux int mu' mu}
\int_n^t   \frac{\mu'(r)}{\mu(r)}\, dr\leq  \ln \mu (n+1)-\ln \mu(n)\leq \ln \theta\leq  \theta,
\end{equation}
we get 
\begin{equation}\label{910x}
\lVert D_x\phi(t, n; x) \rVert \le \hat M 
\end{equation}
for every $t\in [n, n+1]$ and $x\in \R^d$, where $\hat M:=K\theta^ae^{K\eta \theta^{a+1} }$ (which, in particular, is independent of $n$).
On the other hand, note that
\begin{equation}\label{Dfm}
Dg_n(x)=\int_n^{n+1}T(n+1, r)D_xf(r, \phi(r, n;x))D_x\phi(r, n;x)\, dr.
\end{equation}
Then, combining \eqref{bnd}, \eqref{eq: grbound-cont}, \eqref{1951}, \eqref{910x} with \eqref{Dfm} we get
\[
\begin{split}
\lVert Dg_n(x) \rVert &\le \int_n^{n+1}K \bigg (\frac{\mu(n+1)}{\mu(r)} \bigg )^a  \eta \frac{\mu'(r)}{\mu(r)} \hat M  \, dr \\
&\le \frac{K\hat M \theta^{a}\eta }{\mu(n)} \int_n^{n+1}\mu'(r)dr =K\hat M \theta^{a}\eta \left( \frac{\mu(n+1)-\mu(n)}{\mu(n)} \right),
\end{split}
\]
for each $n\in \N$ and $x\in \R^d$. We conclude that \eqref{eq: bound deriv g} holds with $c:=K\hat M\theta^{a}\eta >0$, which can be made sufficiently small by taking $\eta$ sufficiently small. 

Next, we observe that~\eqref{dphi} implies that 
\[
\begin{split}
&D_x\phi(t, n;x)-D_x\phi(t, n; y)
\\
&=\int_n^t T(t,r)D_xf(r, \phi(r, n;x))D_x\phi(r, n;x)\, dr \\
&\phantom{=}-\int_n^t T(t,r)D_xf(r, \phi(r, n;y))D_x\phi(r, n;y)\, dr \\
&=\int_n^t T(t,r)D_xf(r, \phi(r, n;x))(D_x\phi(r, n;x)-D_x\phi(r, n;y))\, dr \\
&\phantom{=}+\int_n^t T(t,r)(D_xf(r, \phi(r, n;x))-D_xf(r, \phi(r, n;y)))D_x\phi(r, n;y)\, dr.
\end{split}
\]
Hence, it follows from~\eqref{bnd}, \eqref{eq: grbound-cont}, \eqref{1951}, \eqref{1952}, \eqref{eq: aux int mu' mu} and~\eqref{910x} (together with the
mean-value theorem) that
\[
\begin{split}
&\lVert D_x\phi(t, n;x)-D_x\phi(t, n; y)  \rVert
\\
&\leq \int_n^t K \bigg (\frac{\mu(t)}{\mu(r)}\bigg )^a L\frac{\mu'(r)}{\mu(r)} \lVert \phi(r, n;x)-\phi(r, n;y)\rVert \hat M \, dr\\
&\phantom{=}+\int_n^t K \bigg (\frac{\mu(t)}{\mu(r)}\bigg )^a \eta \frac{\mu'(r)}{\mu(r)} \lVert D_x\phi(r, n;x)-D_x\phi(r, n;y)\rVert \, dr\\
&\le KL \theta^{a+1}\hat M^2\|x-y\|+ \int_n^t K\theta^a \eta \frac{\mu'(r)}{\mu(r)} \lVert D_x\phi(r, n;x)-D_x\phi(r, n;y)\rVert \, dr,
\end{split}
\]
for $t\in [n, n+1]$ and $x, y\in \R^d$.   By Gronwall's inequality again, one can conclude that there exists $\check M\geq 1$ such that
\begin{equation}\label{Df-Df}
\lVert D_x\phi(t, n;x)-D_x\phi(t, n; y)  \rVert  \le  \check M\lVert x-y\rVert,
\end{equation}
for $n\in \N$, $t\in [n, n+1]$ and $x, y\in \R^d$.

Moreover, since
\[
\begin{split}
&Dg_{n}(x)-Dg_{n}(y)
\\
&=\int_n^{n+1}T(n+1, r)D_xf(r, \phi(r, n;x))D_x\phi(r, n;x)\, dr \\
&\phantom{=}-\int_n^{n+1}T(n+1, r)D_xf(r, \phi(r, n;y))D_x\phi(r, n;y)\, dr \\
&=\int_n^{n+1}T(n+1, r)D_xf(r, \phi(r, n;x))(D_x\phi(r, n;x)-D_x\phi(r, n;y))\, dr \\
&\phantom{=}+\int_n^{n\!+\!1}T(n+1, r)(D_xf(r, \phi(r, n;x))
\!-\!D_xf(r, \phi(r, n;y)))D_x\phi(r, n;y)\, dr, \\
\end{split}
\]
we obtain from \eqref{bnd}, \eqref{eq: grbound-cont}, \eqref{1951}, \eqref{1952}, \eqref{910x} and \eqref{Df-Df} (together with the mean-value
theorem) that
\[
\begin{split}
&\lVert Dg_{n}(x)-Dg_{n}(y)  \rVert
\\
&\le \int_n^{n+1}K \bigg (\frac{\mu(n+1)}{\mu(r)} \bigg )^a \eta \frac{\mu'(r)}{\mu(r)} \check M\lVert x-y\rVert\, dr \\
&\phantom{\le}+\int_n^{n+1}K \bigg (\frac{\mu(n+1)}{\mu(r)} \bigg )^a L\frac{\mu'(r) }{\mu(r)} \hat M  \lVert x-y\rVert \hat M  \, dr \\
&\le \frac{K\check M \theta^{a}\eta}{\mu(n)} \|x-y\| \int_{n}^{n+1}\mu'(r)dr +\frac{K L \hat M^2 \theta^{a} }{\mu(n)} \|x-y\| \int_{n}^{n+1}\mu'(r)dr\\
&= \left(K\check M \theta^{a}\eta + K L \hat M^2 \theta^{a}\right) \frac{\mu(n+1)-\mu(n)}{\mu(n)} \|x-y\|\\
\end{split}
\]
for $n\in \N$ and $x, y\in \R^d$. Then, we conclude that \eqref{4-4-3} holds, and, therefore, all the hypotheses of Theorem \ref{theo: main linearization discrete} are satisfied for $\mathbb{A}=(A_n)_{n\in \N}$ and $(g_n)_{n\in \N}$ defined above. 

Thus, provided that $\eta$ is sufficiently small, by Theorem \ref{theo: main linearization discrete} there exists  a sequence $(\psi_n)_{n\in \N}$ of $C^1$-diffeomorphisms $\psi_n \colon \R^d \to \R^d$ satisfying~\eqref{lincond} and~\eqref{hn-1} (for some $T, \rho >0$).
We now set 
\begin{equation}\label{eq: def of H and G}
    H(t,x)=T(t,n) \psi_n (\phi(n, t;x)) \; \text{ and }\;  G(t,x)=\phi(t, n; \psi_n^{-1}(T(n,t)x)),
\end{equation}
for $n\in \N$, $t\in [n, n+1)$ and $x\in \R^d$. In order to verify that $H(t,x)$ and $G(t,x)$ satisfy properties i) and ii) of our statement, we see from \eqref{eq: def gn cont case} that 
\[
\varphi(n+1,n;x)=A_nx+g_n(x)
\]
for any $n\in \N$ and $x\in \R^d$.
Combining this fact with \eqref{lincond} we get that
\[ 
\begin{split}
H(t,\varphi(t,1;x))&=T(t,n) \psi_n (\phi(n, t;\varphi(t,1;x)))\\
&=T(t,n) \psi_n (\phi(n, 1;x)))\\
&=T(t,n) \psi_n (\phi(n,n-1; \varphi(n-1,1;x)))\\
&=T(t,n) \psi_n ((A_{n-1}+g_{n-1})(\varphi(n-1,1;x)))\\
&=T(t,n) (A_{n-1} \circ \psi_{n-1} (\varphi(n-1,1;x)))\\
&=T(t,n-1) \psi_{n-1} (\varphi(n-1,1;x)),\\
\end{split}
\]
for every $t\in[n,n+1)$ and  $n>1$. The formula above also gives us that 
\[
T(t,n) \psi_n (\phi(n, 1;x)))=
T(t,n-1) \psi_{n-1} (\varphi(n-1,1;x)).
\]
Then, proceeding recursively, we can conclude that 
\[
H(t,\varphi(t,1;x))=T(t,1) \psi_{1} (x) 
\]
for $t\geq 1$. In particular, $H(t,\varphi(t,1;x))$ is a solution of \eqref{LDE} which proves i). Similarly, we can prove that $G$ satisfies property ii) of our statement. 

Now, we observe that
\[
\begin{split}
H(t, G(t, x)) &=T(t, n)\psi_n(\phi(n, t; G(t,x))) \\
&=T(t,n)\psi_n (\phi(n, t; \phi(t, n;\psi_n^{-1}(T(n, t)x)) ) \\
&=T(t,n)\psi_n (\psi_n^{-1}(T(n,t)x)) \\
&=T(t,n)T(n, t)x =x,
\end{split}
\]
for each $x\in \R^d$, $t\in [n, n+1)$ and $n\in \N$, we get $H(t, G(t, x))=x$ for every $t\geq 1$ and $x\in \R^d$. Similarly, we can show that $G(t, H(t, x))=x$ for every $t\geq 1$ and $x\in \R^d$. Consequently, conclusion iii) of our statement also holds true. It remains to show that iv) is satisfied. 

Proceeding as we did to obtain~\eqref{910x}, we find that there exists $\hat M>0$ such that
\begin{equation}\label{Dx-Inv}
\|D_x\phi(n,t;x) \| \le  \hat M
\end{equation}
for $n\in \N$, $t\in [n, n+1)$ and $x\in \R^d$. In particular, recalling that $\varphi(n,t;0)=0$ (which is a consequence of \eqref{1950}) and combining \eqref{Dx-Inv} with the mean-value theorem, we get
\begin{equation}\label{bx1}
\|\phi(n,t;x) \| \le \hat M  \|x\|.
\end{equation}
Therefore, if $\|x\|\leq \frac{ \rho}{\hat M}$ we get that $\|\phi(n,t;x) \| \leq \rho$ for every $t\in[n,n+1)$ and $n\in \N$. Now, since for every
$t\ge 1$ there is an $n\in\mathbb{N}$ such that $t\in [n,n+1)$, combining this observation with \eqref{hn-1}, \eqref{bnd}, \eqref{eq: grbound-cont}, \eqref{Dx-Inv} and the definition of $H$, we get that 
\[
\begin{split}
\|D_xH(t,x)\|&\leq\|T(t,n)\| \|D_x\psi_n (\phi(n, t;x))\|\|D_x\phi(n, t;x)\|\\
&\leq K \theta^a T  \hat M,
\end{split}
\]
whenever $\|x\|\leq \frac{\rho}{\hat M }$. Similarly,
\[
\begin{split}
\|D_xG(t,x)\|&\leq K \theta^a T  \hat M,\\
\end{split}
\]
for every $t\geq 1$ whenever $\|x\|\leq \frac{ \rho}{\hat M }$. Therefore, taking $R=K \theta^a T  \hat M$ and $\zeta =\frac{\rho}{\hat M}$, we obtain iv). This concludes the proof of Theorem \ref{theo: linarization continuous time pol}.
\end{proof}

\end{document}
